\section{Conclusion}
\label{sec:conclusion}

We introduced \eps, a per-token normalized Shannon entropy over top-$k$
log-probabilities, AST-filtered and max-aggregated into a single score
per generated function. \eps is a filter, not a detector: it achieves
the recall target established in \S\ref{subsec:pr}. The contribution is
the two-stage system. Pairing a recall-maximizing filter with a
parallel LLM reviewer produces an architecture in which nothing
uncertain ships unreviewed and the reviewer scales with generation
speed.

The framing matters, and the empirical case for it is stronger than in
earlier iterations of this work. Multi-sample uncertainty methods
produce $10$--$35\%$ detection on our benchmark --- now measured
directly on four models rather than inferred from a single library ---
because code uncertainty is intra-generational: visible in the
single-generation token distribution, invisible across re-samples. The
sequence-level mean token probability \papg, which reads the same raw
logprobs as \eps, is competitive with or better than \eps on raw
detection rate for three of four OpenAI models at the same $1\times$
API cost; \eps's distinguishing property is not higher detection but
token-level attribution, which the stage-2 reviewer uses to target its
analysis. Whether that attribution translates to better end-to-end
precision in a full review loop is an open question we leave for
future work.

The same mechanism --- early-token commitment forcing later tokens ---
plausibly applies in other structured-output domains (SQL,
configuration files, typed program synthesis); empirical verification
is left to future work. Generation is fast; review is the bottleneck.
A filter that misses nothing, plus a parallel reviewer that knows
where to look, moves review off the developer's critical path.
